\chapter{Nomen-Verb-Verbindungen}
\label{cha:NomenVerb}

Viele Nomen kommen in festen Verbindungen mit Verben vor. Diese Verbindungen sind charakteristisch für den deutschen Sprachgebrauch.

\section{Verben mit Akkusativ}
\label{sec:VerbenMitAkk}

\begin{itemize}
    \item \textbf{abgeben}: Ich gebe den Job ab.
    \item \textbf{anfangen}: Wir fangen ein neues Projekt an.
    \item \textbf{anrufen}: Kannst du mich anrufen?
    \item \textbf{aufmachen}: Mach das Fenster auf!
    \item \textbf{ausfüllen}: Ich fülle das Formular aus.
    \item \textbf{besuchen}: Wir besuchen die Großeltern.
    \item \textbf{bezahlen}: Er bezahlt die Rechnung.
    \item \textbf{brauchen}: Ich brauche ein neues Auto.
    \item \textbf{durchsuchen}: Die Polizei durchsucht das Haus.
    \item \textbf{einladen}: Sie lädt ihre Freunde ein.
    \item \textbf{erzählen}: Er erzählt eine Geschichte.
    \item \textbf{finden}: Wir finden die Lösung.
    \item \textbf{fragen}: Ich frage den Lehrer.
    \item \textbf{haben}: Er hat ein Haus.
    \item \textbf{kennen}: Kennst du ihn?
    \item \textbf{kochen}: Sie kocht das Abendessen.
    \item \textbf{legen}: Leg das Buch auf den Tisch!
    \item \textbf{lesen}: Ich lese ein Buch.
    \item \textbf{lieben}: Er liebt seine Frau.
    \item \textbf{mitnehmen}: Nimm das Geschenk mit!
    \item \textbf{mögen}: Ich mag Kaffee.
    \item \textbf{nehmen}: Nimm den Bus!
    \item \textbf{rufen}: Rufen Sie einen Arzt!
    \item \textbf{schenken}: Sie schenkt mir Blumen.
    \item \textbf{schicken}: Ich schicke eine E-Mail.
    \item \textbf{schreiben}: Schreibe einen Brief!
    \item \textbf{sehen}: Ich sehe den Film.
    \item \textbf{setzen}: Setz dich hin!
    \item \textbf{stellen}: Stell die Flasche hin!
    \item \textbf{suchen}: Ich suche meine Brille.
    \item \textbf{tragen}: Er trägt eine Tasche.
    \item \textbf{treffen}: Wir treffen uns um 18 Uhr.
    \item \textbf{trinken}: Trinkst du Kaffee?
    \item \textbf{verstehen}: Verstehst du Deutsch?
    \item \textbf{waschen}: Ich wasche das Auto.
    \item \textbf{werfen}: Wirft du den Ball?
\end{itemize}

\section{Verben mit Dativ}
\label{sec:VerbenMitDat}

Verben, die ein Dativ-Objekt verlangen:

\subsection{Gefühlsverben}
\begin{itemize}
    \item \textbf{gefallen}: Das Buch gefällt mir.
    \item \textbf{gehören}: Das Auto gehört mir.
    \item \textbf{schmecken}: Der Kuchen schmeckt uns.
    \item \textbf{fehlen}: Mir fehlt nur noch ein Buch.
    \item \textbf{passen}: Diese Schuhe passen dir.
\end{itemize}

\subsection{Hilfsverben}
\begin{itemize}
    \item \textbf{helfen}: Ich helfe meiner Schwester.
    \item \textbf{danken}: Ich danke dir.
    \item \textbf{zustimmen}: Ich stimme dir zu.
    \item \textbf{glauben}: Ich glaube dir.
    \item \textbf{folgen}: Der Hund folgt seinem Herrn.
\end{itemize}

\subsection{Objektsverben}
\begin{itemize}
    \item \textbf{antworten}: Er antwortet mir.
    \item \textbf{begegnen}: Wir begegnen ihr.
    \item \textbf{gratulieren}: Ich gratuliere dir.
    \item \textbf{hören}: Ich höre dir zu.
    \item \textbf{zuhören}: Ich höre dir zu.
\end{itemize}

\subsection{Empfindungsverben}
\begin{itemize}
    \item \textbf{wehtun}: Der Arm tut mir weh.
    \item \textbf{leidtun}: Es tut mir leid.
\end{itemize}

\section{Verben mit Dativ und Akkusativ}
\label{sec:VerbenMitDatAkk}

Diese Verben verlangen zwei Objekte:

\begin{itemize}
    \item \textbf{bringen}: Ich bringe meiner Mutter Blumen.
    \item \textbf{empfehlen}: Ich empfehle dir diesen Film.
    \item \textbf{erklären}: Ich erkläre dem Studenten die Regel.
    \item \textbf{geben}: Ich gebe dir das Buch.
    \item \textbf{kaufen}: Ich kaufe meinem Vater eine Uhr.
    \item \textbf{kochen}: Ich koche meinem Mann Suppe.
    \item \textbf{schenken}: Ich schenke meiner Freundin einen Ring.
    \item \textbf{schicken}: Ich schicke dir eine Nachricht.
    \item \textbf{schreiben}: Ich schreibe meinem Chef einen Bericht.
    \item \textbf{stehlen}: Er stiehlt dem Museum ein Gemälde.
    \item \textbf{wünschen}: Ich wünsche dir alles Gute.
    \item \textbf{zeigen}: Ich zeige dem Touristen den Weg.
\end{itemize}

\section{Verben mit Genitiv}
\label{sec:VerbenMitGen}

Diese Verben sind im schriftlichen Deutsch, besonders akademisch, wichtig:

\begin{itemize}
    \item \textbf{bedürfen}: Der Patient bedarf medizinischer Hilfe.
    \item \textbf{gedenken}: Wir gedenken der Toten.
    \item \textbf{sich annehmen}: Sie nahm sich des Problems an.
    \item \textbf{sich bedienen}: Der Professor bedient sich klassischer Zitate.
    \item \textbf{sich erfreuen}: Das Projekt erfreut sich großer Beliebtheit.
    \item \textbf{Herr werden}: Er wurde der Lage Herr.
\end{itemize}

\section{Übungen}
\label{sec:NomenVerbUebungen}

\begin{enumerate}
    \item Kannst du \underline{\hspace{1cm}} helfen? (ich) \quad \textit{(Lösung: mir)}
    \item Das Kleid passt \underline{\hspace{1cm}} sehr gut. (du) \quad \textit{(Lösung: dir)}
    \item Wem glaubst du? (er) \quad \textit{(Lösung: ihm)}
    \item Ich gebe \underline{\hspace{1cm}} das Geschenk. (du) \quad \textit{(Lösung: dir)}
    \item Er zeigt \underline{\hspace{1cm}} das Foto. (wir) \quad \textit{(Lösung: uns)}
\end{enumerate}

\chapterend